% !PRINTED none
% Chapter-opening page: centred "Zweiter Abschnitt." / "Theta-Funktionen." with a rule
% beneath, then the \S\,15 heading and running text. No page number is printed on
% this page (chapter-opening convention).

\begin{center}
\gesperrt{Zweiter Abschnitt.}

\medskip
\textbf{\gesperrt{Theta-Funktionen.}}
\end{center}

\bigskip
\begin{center}\rule{2cm}{0.4pt}\end{center}
\bigskip

\begin{center}
\textbf{\S\,15. \gesperrt{Voraussetzungen aus der Funktionentheorie.}}
\end{center}

Der Begriff der doppelt periodischen Funktionen kann nur dann richtig aufgefaßt werden, wenn man sie als Funktionen eines komplexen Arguments $u = v + iw$ betrachtet, worin $i$ wie gewöhnlich die Bedeutung von $\sqrt{-1}$ hat. Die Variable $u$ gilt uns als unabhängige und unbeschränkt veränderliche Größe, die nach dem Vorgange von \gesperrt{Gauss} durch die Punkte einer Ebene in der Weise geometrisch veranschaulicht wird, daß der Punkt, dessen Koordinaten in einem rechtwinkeligen System $v, w$ sind, als Träger des Wertes $u = v + iw$ angesehen und kurz als der Punkt $u$ bezeichnet wird.

Wir beschränken unsere Betrachtung hier auf eindeutige analytische Funktionen $\varphi(u)$, die im Endlichen keine \gesperrt{wesentlich singulären} Stellen haben. Was im Unendlichen geschieht, lassen wir dahingestellt. Wir setzen also von den betrachteten Funktionen folgendes voraus:

1. In jedem endlichen Flächenstück liegt eine endliche Anzahl von Punkten, in denen die Funktion $\varphi(u)$ unendlich wird; diese Punkte heißen \gesperrt{Unstetigkeitspunkte}.

Die Anzahl der Unstetigkeitspunkte überhaupt, d.\,h. in der ganzen unendlichen Ebene, kann natürlich auch unendlich groß sein.

2. Ist $u_0$ ein nicht zu den Unstetigkeitspunkten gehöriger Punkt, so ist die Funktion entwickelbar in eine nach ganzen aufsteigenden positiven Potenzen von $u - u_0$ fortschreitende Reihe, die konvergent ist in einem Kreise, der den Punkt $u_0$ zum Mittelpunkte hat und bis zum nächstgelegenen Unstetigkeitspunkte reicht. Man sagt, die Funktion habe in der Umgebung des Punktes $u_0$ den \gesperrt{Charakter einer ganzen Funktion}.
